\documentclass[11pt]{article}

\usepackage[margin=1in]{geometry}
\usepackage{booktabs}
\usepackage{longtable}
\usepackage{xcolor}
\usepackage{hyperref}
\usepackage{amsmath}
\usepackage{enumitem}
\usepackage{titlesec}
\usepackage{fancyhdr}

\definecolor{verdictred}{RGB}{150,30,30}
\definecolor{headergray}{RGB}{90,90,90}
\definecolor{boxgray}{RGB}{245,245,245}

\hypersetup{colorlinks=true, linkcolor=black, citecolor=black,
  urlcolor={blue!60!black},
  pdftitle={PEARL-Neuro EEG Classification --- Final Report},
  pdfauthor={PEARL-Neuro Analysis Team}}

\titleformat{\section}{\large\bfseries}{\thesection}{0.6em}{}
\titleformat{\subsection}{\normalsize\bfseries}{\thesubsection}{0.6em}{}

\pagestyle{fancy}
\fancyhf{}
\fancyhead[L]{\small\color{headergray} PEARL-Neuro --- Final Report (v2)}
\fancyhead[R]{\small\color{headergray} \thepage}
\renewcommand{\headrulewidth}{0.4pt}

\newcommand{\keybox}[1]{%
  \begin{center}\fcolorbox{black}{boxgray}{\parbox{0.94\linewidth}{\vspace{2pt}#1\vspace{2pt}}}\end{center}}

\title{\vspace{-1.5cm}PEARL-Neuro EEG Classification\\[4pt]
\large Predicting Genetic Alzheimer's-Risk Group from Resting-State EEG\\[2pt]
\large Final Report --- Version 2 (consolidated)}
\author{Dataset: OpenNeuro \texttt{ds004796} --- 79 genotyped, EEG-imaged subjects}
\date{2026-08-15}

\begin{document}
\maketitle
\thispagestyle{fancy}

\keybox{\textbf{\color{verdictred}Headline.} Resting-state EEG did not distinguish
the genetic Alzheimer's-risk groups in this cohort. This is a \emph{bounded}
null, not an inconclusive one: the same features that find nothing for genotype
successfully detect a real subject-level trait (sex) on the identical cohort and
analysis protocol. The measurement works; the signal is not there.}

\tableofcontents

% =====================================================================
\section{What was asked, what was delivered}

\textbf{What was asked.} Build a classifier that predicts genetic
Alzheimer's-risk group (derived from APOE~$\varepsilon$4 and PICALM genotype)
from EEG, using a pitch-synchronous waveform-shape feature family (PSR/PSWT)
compared against a conventional spectral baseline, benchmarked against a
published result of AUC~$\approx$~0.58.

\textbf{What was delivered.} A validated, reproducible, pre-registered EEG
analysis pipeline, and the bounded null result it produced --- together with a
quantified statement of exactly how large an effect this dataset was capable of
detecting. No usable predictor was delivered, because none exists to deliver:
the target is not detectable in this cohort at this sample size, and that
conclusion is now supported by direct evidence rather than assumed from absence.

% =====================================================================
\section{The headline, in plain language}

Think of the pipeline as a metal detector and the dataset as a field. Sweeping
the field found no gold. That alone is ambiguous --- it could mean the field is
empty, or that the detector is broken.

We buried a known coin and confirmed the detector beeps. Using the identical
features, subjects, preprocessing and statistics, the pipeline reliably detects a
subject's sex (a trait EEG genuinely encodes) while detecting nothing at all for
genotype group.

\keybox{The field really is empty. The genotype null is a statement about
genotype, not about the sensitivity of the measurement.}

% =====================================================================
\section{The three acts}

\subsection{Act I --- the original null}

The pre-registered primary analysis returned \textbf{AUC 0.474}
(95\% CI [0.266, 0.702], permutation $p = 0.577$) on $n = 64$ subjects: no
better than chance.

The pipeline was built so that a positive result would have been believable ---
label-blind preprocessing enforced by tests, a cross-validation and inference
protocol frozen before any label was joined, and provenance tracking from raw
file hash to final number.

\subsection{Act II --- the near-miss, and its retraction}

Reviewing the pipeline surfaced two genuine preprocessing defects: the average
EEG reference was computed \emph{before} bad-channel interpolation (letting one
bad channel contaminate every other channel), and the ICA component count was
fixed rather than scaled per subject. Both were fixed on their own merits.

Re-running the \emph{identical} frozen analysis on the corrected data flipped the
verdict:

\begin{table}[h]\centering
\begin{tabular}{lcc}
\toprule
 & Original & After the two fixes \\
\midrule
AUC & 0.474 & \textbf{0.651} \\
95\% CI & [0.266, 0.702] & [0.353, 0.764] \\
Permutation $p$ & 0.577 & \textbf{0.041} \\
Verdict & null & \emph{apparently} positive \\
\bottomrule
\end{tabular}
\end{table}

This looked like a real finding. It exceeded the published 0.58 benchmark and
cleared conventional significance. \textbf{It was then subjected to three
independent checks and did not survive any of them:}

\begin{itemize}[leftmargin=1.2cm]
  \item \textbf{Stability.} Resampling the same 64 subjects put the median AUC at
        \textbf{0.565} --- essentially on the covariates-only baseline (0.557) ---
        with \textbf{24\%} of resamples falling below chance entirely.
  \item \textbf{Attribution.} Isolating each preprocessing fix in a separate full
        re-run showed \emph{neither} reproduces significance alone: reference-order
        fix 0.583 ($p = 0.234$), ICA fix 0.602 ($p = 0.174$). Only the specific
        combination, on this specific sample, crossed $p < 0.05$.
  \item \textbf{Reproducibility.} The 0.651 figure regenerated exactly (0.6509,
        $p = 0.0410$) at full precision, confirming it was computed correctly ---
        but correctness of computation is not evidence of a real effect.
\end{itemize}

The result was retracted. It is reported here in full because concealing it would
misrepresent what is known.

\subsection{Act III --- bounding the null}

One ambiguity remained. An early diagnostic had shown the PSWT feature family
failing a basic sanity check: it could not detect a subject's \emph{sex}
(AUC 0.465 against a pre-declared 0.65 floor), even though resting EEG normally
encodes sex. So ``we found nothing'' might have meant ``these features discard
all between-subject information.''

Four representations already computed from the corrected data were tested against
the sex control --- no new analysis of the genotype target was run:

\begin{table}[h]\centering
\caption{Sex control, $n=64$, full precision (1000 permutations / 2000 bootstrap).}
\begin{tabular}{lccccc}
\toprule
Representation & Cols & AUC & 95\% CI & perm $p$ & Verdict \\
\midrule
\texttt{baseline\_rest} & 9 & \textbf{0.688} & [0.431, 0.816] & \textbf{0.019} & PASS \\
\texttt{baseline\_rest\_msit} & 16 & \textbf{0.696} & [0.421, 0.826] & \textbf{0.007} & PASS \\
\texttt{baseline\_all} & 23 & \textbf{0.652} & [0.373, 0.819] & \textbf{0.029} & PASS \\
\texttt{pswt} & 18 & 0.465 & [0.272, 0.700] & 0.619 & FAIL \\
\bottomrule
\end{tabular}
\end{table}

Three of four decode sex significantly; only the pitch-synchronous family fails.
The decisive comparison holds the representation fixed and changes only the
target --- the same nine columns:

\begin{table}[h]\centering
\begin{tabular}{lccl}
\toprule
Target & AUC & $p$ & Outcome \\
\midrule
Sex (control) & \textbf{0.688} & \textbf{0.019} & signal detected \\
Genotype (primary) & 0.506 & 0.482 & nothing \\
\bottomrule
\end{tabular}
\end{table}

Matched on features, cohort, preprocessing and protocol. \textbf{The null is
bounded.}

% =====================================================================
\section{Why the retracted result is the strongest evidence of rigour here}

This deserves stating plainly, because it is easy to read Act II as a setback.

A team that ran the corrected preprocessing, obtained AUC 0.651 at $p = 0.041$,
and stopped there would have reported a finding that beat the published
benchmark. It would have been published, or acted upon, or built into a product.
\textbf{And it would have been false} --- as the resampling and attribution
checks subsequently demonstrated.

The value of this engagement is not the answer it produced. It is that the
process was built to kill its own most attractive result when that result failed
validation, and did so. A pipeline that only ever confirms is worthless; this one
demonstrably does not.

% =====================================================================
\section{The detection ceiling --- a finding about the dataset}

A formal power analysis ($n = 34$ vs 30, $\alpha = 0.05$ two-sided) quantifies
what this cohort could ever have resolved:

\begin{table}[h]\centering
\begin{tabular}{lc}
\toprule
Quantity & Value \\
\midrule
Minimum detectable AUC at 80\% power & \textbf{0.693} \\
Minimum detectable AUC at 90\% power & 0.720 \\
Power to detect the published benchmark (0.58) & \textbf{20\%} \\
Power to detect AUC 0.65 & 57\% \\
\bottomrule
\end{tabular}
\end{table}

Two consequences worth carrying forward:

\begin{itemize}[leftmargin=1.2cm]
  \item \textbf{All three passing sex controls sit at or near that 0.693 floor.}
        Sex is among the most robustly encoded subject-level traits in resting
        EEG, and this cohort can only just resolve it. A subtler effect has no
        headroom.
  \item \textbf{At 20\% power against the benchmark effect}, a study of this size
        would fail to detect that effect four times in five \emph{even if it were
        real and exactly as large as published}.
  \item \textbf{Age is unusable as a control in this cohort} --- it spans only
        50--63 years, and published EEG brain-age error ($\approx$ 7 years)
        exceeds that entire range.
\end{itemize}

These are findings \emph{about the dataset} and are useful to anyone else
attempting this question on it.

% =====================================================================
\section{What was actually built}

\keybox{\textbf{The deliverable is a validated, reproducible EEG analysis
pipeline and the bounded null it produced --- not a predictor.}}

Concretely: BIDS ingestion $\rightarrow$ label-blind preprocessing (including the
two fixes, now permanent defaults) $\rightarrow$ PSR/PSWT and spectral feature
extraction $\rightarrow$ a pre-registered cross-validation, permutation and
bootstrap evaluation harness $\rightarrow$ provenance-traced reporting.

Properties that make it reusable rather than disposable:

\begin{itemize}[leftmargin=1.2cm]
  \item \textbf{Label blindness enforced by tests}, not convention --- the
        preprocessing and feature packages have no import path to the genotype
        labels, asserted by dedicated tests.
  \item \textbf{Pre-registration.} The analysis plan, CV scheme, model class and
        stopping rules were frozen before any label was joined; changes are
        append-only amendments.
  \item \textbf{Provenance.} Every derivative records the source file hash, the
        processing spec, a run ID and the git commit that produced it.
  \item \textbf{289 passing tests}, including real-data regression checks.
\end{itemize}

A future, adequately powered cohort would be run through this pipeline unchanged.

% =====================================================================
\section{What this does and does not license}

\textbf{This report does \emph{not} say} that APOE/PICALM status is invisible in
EEG, or that the question is closed in general.

\textbf{It says} that these representations, on this cohort, at this sample size,
did not detect it --- and it quantifies how large an effect would have had to be
to show up (AUC $\geq$ 0.693 for 80\% power). An effect smaller than that could
be entirely real and would have been invisible here.

The distinction matters for anyone deciding whether to pursue this line further:
the evidence bounds the effect, it does not exclude it.

% =====================================================================
\section{What a definitive answer would require}

\begin{enumerate}[leftmargin=1.2cm]
  \item \textbf{A new, adequately powered cohort}, sized from the power analysis
        on record rather than from availability. This is the binding constraint;
        no methodological change substitutes for it.
  \item \textbf{Independent replication data.} \texttt{ds004796} appears to be
        the only public EEG dataset genotyped for APOE/PICALM at this scale, so
        replication requires new primary collection.
  \item \textbf{Task-state EEG}, if the published literature's claim that task
        outperforms rest for this question is to be tested properly --- this
        project's own task-based benchmark reproduction landed at chance and the
        gap remains characterised but unresolved.
\end{enumerate}

% =====================================================================
\section{Limitations and known defects}

\subsection{Limitations}

\begin{itemize}[leftmargin=1.2cm]
  \item \textbf{Hard sample ceiling.} 79 subjects have both EEG and genotype; all
        were used. 64 survive QC for the rest-based primary analysis.
  \item \textbf{No independent replication cohort exists} (see above).
  \item \textbf{Benchmark mismatch.} The published 0.58 came from task-state
        (MSIT) EEG with a substantially richer feature pipeline on a different
        cohort; it is a characterised methodological difference, not a validated
        comparison point in either direction.
  \item \textbf{Cohort sensitivity.} A change of only 5 subjects (a QC side
        effect) moved the point estimate by 0.177 AUC --- inherent to this sample
        size.
  \item \textbf{The bound applies to representations tested.} Three independent
        representations carry subject-level information and find no genotype
        effect; this does not prove no representation could.
\end{itemize}

\subsection{Known defects annex}

Recorded rather than concealed:

\begin{itemize}[leftmargin=1.2cm]
  \item \textbf{Report-writer overwrite defect (fixed).} Report writers used
        fixed output paths with no overwrite protection. This silently destroyed
        131 lines of the pre-registered analysis plan (two appended amendments)
        during re-runs; recovered from version control. Report writes are now
        guarded --- the frozen plan is write-once and all other reports archive
        any prior version before overwriting.
  \item \textbf{Hardcoded narrative (fixed).} Two report generators contained
        prose describing a specific expected result, including a stale
        ``verdict (frozen: NULL, AUC 0.474)'' claim that would have been emitted
        unchanged alongside a different result. Narrative is now generated from
        computed values or not at all, enforced by test.
  \item \textbf{Duplicated disclaimer (fixed).} The inference module served by
        the container carried its own copy of the output disclaimer, so
        strengthening the primary definition would have left the shipped surface
        unchanged. Now a single imported definition.
  \item \textbf{Control-figure provenance (corrected).} An earlier report draft
        described the sex control --- a \emph{failed positive} control --- as a
        passed negative control, and did not state which preprocessing each
        control figure came from. Corrected: 0.849 / 0.465 / 0.490 are from the
        corrected preprocessing; 0.830 / 0.486 / 0.459 from the original.
\end{itemize}

% =====================================================================
\section{Deliverables index}

\begin{longtable}{ll}
\toprule
Artefact & Path \\
\midrule
\endhead
This report & \texttt{reports/final\_report\_v2.tex} \\
Bounded null (Phase 6) & \texttt{reports/phase6\_bounded\_null.md} \\
Detection ceiling & \texttt{reports/phase6\_power\_analysis.txt} \\
Retraction of the 0.651 result & \texttt{reports/phase5\_stage3\_validation.md} \\
Stop decision (no further analysis) & \texttt{reports/phase7\_stop\_decision.md} \\
Delivery decision (what ships) & \texttt{reports/phase7\_delivery\_decision.md} \\
Pre-registered analysis plan & \texttt{reports/analysis\_plan\_frozen.md} \\
Original Phase 3 model (frozen) & \texttt{data/derivatives/models\_phase3to4\_frozen/} \\
Reference model artefact + card & \texttt{data/derivatives/models\_stage3/} \\
Reproduction instructions & \texttt{RUNBOOK.md}, \texttt{INSTALL.md} \\
\bottomrule
\end{longtable}

\vspace{0.4cm}
\hrule
\vspace{0.3cm}
{\small\color{headergray}\noindent
This project answered its question. The answer was ``no''. The reason that answer
is worth having is that the pipeline which produced it was built so a ``yes''
would have been believable --- pre-registered, label-blind, provenance-traced,
and willing to discard its own best result when that result failed validation.}

\end{document}
